% Part: lambda-calculus
% Chapter: syntax
% Section: de-bruijn

\documentclass[../../../include/open-logic-section]{subfiles}

\begin{document}

\olfileid{lam}{syn}{deb}

\olsection{দ্য ব্রুইন সূচক}

$\alpha$-সমতুল্যতা খুবই স্বাভাবিক ধারণা, কারণ $\alpha$-সমতুল্য পদগুলির
``অর্থ একই''। বস্তুত, ল্যাম্বডা পদের এমন বাক্যগঠন দেওয়া সম্ভব, যা
$\alpha$-রূপান্তরের মাধ্যমে পরস্পরে বদলানো যায় এমন পদগুলির মধ্যে পার্থক্য
করে না। এর সবচেয়ে পরিচিত রূপে চলরাশির জায়গায় তাদের
\emph{দ্য ব্রুইন সূচক} বসানো হয়।

আমরা যখন $\lambd[x][M]$ লিখি, তখন স্পষ্টভাবে জানাই যে~$x$ অপেক্ষকটির
প্যারামিটার; ফলে~$M$-এ এই প্যারামিটারকে নির্দেশ করতে~$x$ ব্যবহার করা
যায়। কিন্তু দ্য ব্রুইন সূচকে প্যারামিটারের কোনো নাম থাকে না; অপেক্ষকের
দেহে তাকে এমন একটি সংখ্যা দিয়ে নির্দেশ করা হয়, যা তার ও সংশ্লিষ্ট
বিমূর্তনের মধ্যবর্তী বিমূর্তনের স্তরসংখ্যা জানায়। যেমন,
$\lambd[x][\lambd[y][y x]]$ পদটি বিবেচনা করো: বাইরের বিমূর্তনটি
চলরাশি~$x$-কে আবদ্ধ করে, আর ভিতরের বিমূর্তনটি চলরাশি~$y$-কে আবদ্ধ করে।
উপপদ~$y x$ ভিতরের বিমূর্তনের পরিসরে; $y$ ও তাকে আবদ্ধকারী
$\lambd[y]$-এর মধ্যে কোনো বিমূর্তন নেই, কিন্তু~$x$ ও তাকে আবদ্ধকারী
$\lambd[x]$-এর মধ্যে একটি বিমূর্তন আছে। তাই $y x$-এর জন্য আমরা
$0\, 1$, আর সমগ্র পদের জন্য $\lambd[][\lambd[][01]]$ লিখি।

\begin{defn}
  দ্য ব্রুইন পদগুলি নিচের নিয়মে আরোহীভাবে সংজ্ঞায়িত:
  \begin{enumerate}
  \item $n$, যেখানে $n$ যেকোনো স্বাভাবিক সংখ্যা।
  \item $PQ$, যেখানে $P$ ও~$Q$ উভয়ই দ্য ব্রুইন পদ।
  \item $\lambd[][N]$, যেখানে $N$ একটি দ্য ব্রুইন পদ।
  \end{enumerate}
\end{defn}

সাধারণ ল্যাম্বডা পদ থেকে দ্য ব্রুইন সূচকযুক্ত পদে একটি আনুষ্ঠানিক অনুবাদ
নিচে দেওয়া হলো:
\begin{defn}
  \begin{align*}
    F_\Gamma(x) &= \Gamma(x) \\
    F_\Gamma(PQ) &= F_\Gamma(P)F_\Gamma(Q) \\
    F_\Gamma(\lambd[x][N]) &= \lambd[][F_{x,\Gamma}(N)]
  \end{align*}
  এখানে $\Gamma$ হলো শূন্য থেকে সূচকিত চলরাশির একটি তালিকা, এবং
  $\Gamma(x)$ হলো $\Gamma$-তে চলরাশি~$x$-এর অবস্থান। যেমন, $\Gamma$
  যদি $x,y,z$ হয়, তবে $\Gamma(x)$ হলো~$0$ এবং $\Gamma(z)$ হলো~$2$।

  $x,\Gamma$ দিয়ে $\Gamma$-র শুরুতে~$x$ বসিয়ে পাওয়া তালিকা বোঝানো
  হয়। যেমন, আগের উদাহরণের $\Gamma$ ধরে $w,\Gamma$ হলো $w,x,y,z$।
\end{defn}

দ্য ব্রুইন পদ থেকে একটি সাধারণ ল্যাম্বডা পদ নিচের নিয়মে পুনরুদ্ধার করা
হয়:

\begin{defn}
  \begin{align*}
    G_\Gamma(n) &= \Gamma[n] \\
    G_\Gamma(PQ) &= G_\Gamma(P) G_\Gamma(Q) \\
    G_\Gamma(\lambd[][N]) &= \lambd[x][G_{x,\Gamma}(N)]
  \end{align*}
  এখানে $\Gamma$ আবারও শূন্য থেকে সূচকিত চলরাশির একটি তালিকা, এবং
  $\Gamma[n]$ হলো $n$-তম অবস্থানের চলরাশি। যেমন, $\Gamma$ যদি $x,y,z$
  হয়, তবে $\Gamma[1]$ হলো~$y$।

  শেষ সমীকরণের চলরাশি~$x$ হিসেবে $\Gamma$-তে নেই এমন যেকোনো চলরাশি
  বেছে নেওয়া হয়।
\end{defn}

এখানে কয়েকটি ফল প্রমাণ ছাড়াই দেওয়া হলো:

\begin{prop}
  $M \aconvone M'$ এবং $\Gamma$ যদি $\FV{M}$-কে ধারণকারী যেকোনো তালিকা
  হয়, তবে $F_\Gamma(M) \eqs F_\Gamma(M')$।
\end{prop}

\end{document}
